\doublespacing % Do not change - required

\chapter{Literature Review}
\label{ch:literature-review}

%%%%%%%%%%%%%%%%%%%%%%%%%%%%%%%%%%%%%%%
% IMPORTANT
\begin{spacing}{1} %THESE FOUR
\minitoc % LINES MUST APPEAR IN
\end{spacing} % EVERY
\thesisspacing % CHAPTER
% COPY THEM IN ANY NEW CHAPTER
%%%%%%%%%%%%%%%%%%%%%%%%%%%%%%%%%%%%%%%

\section{Purpose and scope}

Use this chapter to establish the thesis-level research context before the paper-based chapters begin. The literature review should explain the main bodies of work that the thesis builds on, compare their assumptions and limitations, and identify the gap that motivates the four publications.

For a thesis by publication, this chapter should not simply repeat the related work sections from the papers. Instead, it should synthesise the shared intellectual background across the thesis and make clear how the later paper chapters fit within that landscape.

\section{Research context and core concepts}

Define the key research area, concepts, and terminology used across the thesis. This section should give readers enough background to understand the later paper-based chapters without repeating the full introduction of each publication.

Organise the discussion around the conceptual or technical distinctions that matter for the thesis. Avoid presenting the literature as a list of papers; group work by the assumptions, mechanisms, or limitations that shape the research gap.

\section{Main approaches in prior work}

Summarise the main approaches, theories, methods, datasets, systems, or analytical traditions that prior work has used to address the problem. For each group of work, explain the common paradigm before discussing representative examples.

End each subsection or paragraph by identifying the limitation that is relevant to the thesis. The limitation should be specific enough to motivate the later contributions rather than a generic statement that more research is needed.

\section{Closest related work and thesis gap}

Discuss the work most closely related to the thesis contributions. This section should make the novelty of the thesis easy to verify by comparing mechanisms, assumptions, evidence, and failure modes.

Clarify how the thesis differs from the closest prior work in precise terms. Depending on the field, this may involve a different research question, method, dataset, theory, system design, empirical setting, evaluation protocol, or interpretation of results.

\section{Evaluation practices and practical constraints}

Summarise the evaluation practices used in the relevant literature, including common evidence types, datasets, corpora, field sites, metrics, experimental protocols, theoretical criteria, or qualitative methods. This section should make clear which aspects of the thesis can be compared directly to prior work and which require a different evaluation framing.

Use this section to surface practical constraints such as data availability, measurement validity, annotation or recruitment cost, reproducibility, ethical considerations, computational cost, deployment assumptions, or transfer across settings.

\section{Summary and thesis positioning}

Conclude the literature review by restating the gap that the thesis addresses. This section should act as a bridge into Chapter~\ref{ch:paper1}: explain why the first paper-based chapter is the natural first technical step after the literature review.

The final paragraph should make the chapter sequence feel deliberate: Chapter~\ref{ch:paper1} introduces the first contribution, Chapters~\ref{ch:paper2}--\ref{ch:paper4} extend or evaluate the thesis from complementary angles, and Chapter~\ref{ch:synthesis} later integrates the findings.
